\documentclass[11pt]{letter}
\usepackage[margin=1in]{geometry}
\usepackage{hyperref}

\signature{Tejas R \\ tejasr.me@rvce.edu.in \\ Department of Mechanical Engineering \\ RV College of Engineering}

\begin{document}

\begin{letter}{
    Editor-in-Chief \\
    Advances in Space Research (ASR) \\
    Elsevier
}

\opening{Dear Editor,}

We are pleased to submit our manuscript, "Discrete Element Modeling of Lunar Regolith Tribocharging in Vibratory Feeders under Reduced Gravity using LIGGGHTS," for consideration as an original research article in \textit{Advances in Space Research}.

This study addresses a critical bottleneck in the In-Situ Resource Utilization (ISRU) value chain: the triboelectric beneficiation of lunar regolith. While electrostatic separation is a leading candidate for mineral extraction on the Moon, the effectiveness of these systems is fundamentally governed by the mechanical feed mechanisms that induce differential tribocharge. Our research demonstrates that under reduced lunar gravity (1.62 m/s$^2$), passive feed geometries fail due to a gravity-dependent shift in the Granular Bond Number ($Bo_g$), which transitions the regolith into a highly cohesive, jammed state.

To investigate this, we developed a coupled computational methodology integrating a Discrete Element Method (DEM) framework with an energy-based tribocharging post-processor. By strictly calibrating the Hertz-Mindlin contact model with SJKR cohesion against the physical angle of repose of the LMS-1 lunar simulant, we ensure high physical fidelity in our simulations. Our parametric analysis quantifies that lunar feeders require significantly higher vibration amplitudes than their terrestrial counterparts to overcome cohesion and achieve the minimum required charge-to-mass ratio ($Q/M$) for effective beneficiation.

We believe this manuscript aligns perfectly with the scope of \textit{Advances in Space Research}, as it provides quantitative design guidelines bridging computational granular physics with applied space systems engineering. 

This manuscript is original, has not been published previously, and is not under consideration for publication elsewhere. All authors have read and approved the final manuscript.

\vspace{1em}
\textbf{Suggested Reviewers:}
\begin{enumerate}
    \item Dr. Carlos I. Calle (\texttt{carlos.i.calle@nasa.gov}) - Lead of the Electrostatics and Surface Physics Laboratory, NASA Kennedy Space Center. Expert in lunar electrostatic beneficiation, triboelectric charging, and space dust physics.
    \item Dr. Philip T. Metzger (\texttt{philip.metzger@ucf.edu}) - Planetary Scientist at the Florida Space Institute, University of Central Florida. Expert in DEM modeling of planetary regolith, granular dynamics under reduced gravity, and lunar simulants.
    \item Dr. Siegfried Ku{\ss} (\texttt{siegfried.kuss@dlr.de}) - Research Engineer at the German Aerospace Center (DLR). Expert in dry electrostatic separation testbed design and regolith feed mechanisms.
\end{enumerate}

Thank you for your time and consideration. We look forward to your response.

\closing{Sincerely,}

\end{letter}
\end{document}
